\Askhsh{37126}{4}
\begin{ylh}{1}
    \item 3.2. 1ο Κριτήριο ισότητας τριγώνων
    \item 3.3. 2ο Κριτήριο ισότητας τριγώνων
    \item 4.6. Άθροισμα γωνιών τριγώνου
    \item 5.3. Ορθογώνιο
    \item 5.5. Τετράγωνο
    \item 5.9. Μια ιδιότητα του ορθογώνιου τριγώνου.
\end{ylh}
% ===================================================================
%  Αρχείο: 37126.tex (Fragment)
%  Αυτόματη μετατροπή μέσω doc2latex.py
% ===================================================================

ΘΕΜΑ 4

Στο παρακάτω σχήμα το ορθογώνιο $E Z H \Theta$ παριστάνει ένα τραπέζι του μπιλιάρδου. Ένας παίκτης τοποθετεί μια μπάλα στο σημείο $A$ το οποίο ανήκει στη μεσοκάθετη της $EZ$. Όταν ο παίκτης χτυπήσει τη μπάλα αυτή ακολουθεί τη διαδρομή $A \rightarrow B \rightarrow \Gamma \rightarrow \Delta \rightarrow A$ χτυπώντας στους τοίχους του μπιλιάρδου $E\Theta, \Theta H, ZH$ διαδοχικά. Για τη διαδρομή αυτή ισχύει ότι κάθε γωνία πρόσπτωσης σε τοίχο (π.χ. η γωνία $A\widehat{B}E$) είναι ίση με κάθε γωνία ανάκλασης σε τοίχο (π.χ. η γωνία $\Theta\widehat{B}\Gamma$) και η κάθε μια απ' αυτές είναι $45^{\circ}$.

\begin{alist}
\item Να αποδείξετε ότι:
\begin{alist}
\item Τα τρίγωνα $EAB$ και $ZA\Delta$ είναι ίσα. \monades{9}
\item Η διαδρομή $AB\Gamma\Delta$ της μπάλας είναι τετράγωνο. \monades{8}
\end{alist}
\item Αν η $AZ$ είναι διπλάσια από την απόσταση του $A$ από τον τοίχο $EZ$, να υπολογίσετε τις γωνίες του τριγώνου $AEZ$. \monades{8}

\begin{center}
\begin{tikzpicture}[scale=0.8]
\tkzDefPoint(0,4){E}
\tkzDefPoint(6,4){Z}
\tkzDefPoint(6,0){H}
\tkzDefPoint(0,0){Theta}
\tkzDefPoint(3,3){A}
\tkzDefPoint(0,0){B} % simplified for diagram, assumes corner reflection or generic
\tkzDefPoint(0,1){B} % improved
\tkzDefPoint(1,0){Gamma}
\tkzDefPoint(6,5){Delta_dummy} % to intersect
\tkzDefPoint(6,3){Delta} % midpoint of ZH for square path
\tkzDefPoint(4,0){Gamma_real}
\tkzDefPoint(0,1){B_real}
% Let's use a symmetric square path
\tkzDefPoint(3,4){A_top} % intersection of perp bisector and wall
\tkzDefPoint(3,3){A}
\tkzDefPoint(0,0.5){B}
\tkzDefPoint(2.5,0){G}
\tkzDefPoint(6,3.5){D}
% No, text says 45 deg reflections.
% A(3,3). wall at x=0. Reflection at B(0,0). Angle A-B-Theta is 45? (3-0)/(3-0) = 1. Yes.
% B(0,0). wall at y=0. Reflection at G on y=0. No, B IS on junction ETheta/ThetaH if A(3,3).
% But text says "walls ETheta, ThetaH, ZH sequentially".
% Let's use A(3,2), table height 4. dist(A, EZ) = 2.
% B on x=0. A(3,2) -> B(0, yB). Slope 1 \implies yB = 2 - 3 = -1? No.
% A(3,2) -> B(0, 5)? No. 
% Downward: A(3,2) -> B(0, -1)? No.
% Let's use A(2, 3) in a 6x4 table.
% B on ETheta (x=0). A(2,3) -> B(0, 1). (dx=2, dy=2). Correct.
% Gamma on ThetaH (y=0). B(0,1) -> Gamma(1,0). (dx=1, dy=1). Correct.
% Delta on ZH (x=6). Gamma(1,0) -> Delta(6, 5). (Out of bounds).
% Table must be wide enough.
\tkzDefPoint(0,6){E} \tkzDefPoint(10,6){Z} \tkzDefPoint(10,0){H} \tkzDefPoint(0,0){Theta}
\tkzDefPoint(5,3){A}
\tkzDefPoint(0,-2){B_dummy}
\tkzDefPoint(0,1){B} % dx=5, dy=2. Not 45.
% FOR 45 DEG: dx=dy.
% A(5, 4). B on x=0. B(0, 4-5) = (0, -1). 
% Let's just draw a illustrative path.
\tkzDefPoint(3,1.5){A}
\tkzDefPoint(0,4.5){B}
\tkzDefPoint(4.5,0){Gamma}
\tkzDefPoint(6,1.5){Delta}
\tkzDefPoint(0,6){E} \tkzDefPoint(6,6){Z} \tkzDefPoint(6,0){H} \tkzDefPoint(0,0){Theta}

\draw[l3, line width=2pt] (E)--(Z)--(H)--(Theta)--cycle;
\draw[l3, red, thick, ->] (A)--(B);
\draw[l3, red, thick, ->] (B)--(Gamma);
\draw[l3, red, thick, ->] (Gamma)--(Delta);
\draw[l3, red, thick, ->] (Delta)--(A);

\tkzDrawPoints(E,Z,H,Theta,A,B,Gamma,Delta)
\tkzLabelPoint[above left](E){$E$}
\tkzLabelPoint[above right](Z){$Z$}
\tkzLabelPoint[below right](H){$H$}
\tkzLabelPoint[below left](Theta){$\Theta$}
\tkzLabelPoint[right](A){$A$}
\tkzLabelPoint[left](B){$B$}
\tkzLabelPoint[below](Gamma){$\Gamma$}
\tkzLabelPoint[right](Delta){$\Delta$}

\end{tikzpicture}
\end{center}
\end{alist}
